\chapter{Requirement Analysis}
\label{chap:requirement-analysis}

\section{Stakeholder Analysis}
\label{section:stakeholder-analysis}

The primary stakeholder of SantaTrail is the beginner drone pilot, a user with little or no prior experience in UAV mission planning who needs a safe and engaging learning environment.

In SantaTrail, the system is designed to teach mission planning concepts in an engaging and progressive manner. The user completes a series of present delivery missions that gradually increase in difficulty, allowing them to build skills over time. The system supports this learning process through gamified, mission-based training that is integrated into a cohesive Christmas narrative, creating a more immersive experience. Additionally, the platform incorporates realistic UAV simulation by connecting missions to QGroundControl and ArduPilot SITL via MAVLink, enabling users to practice in an environment that closely mirrors real-world drone operations. To further enhance learning, AI-based adaptive training dynamically adjusts mission difficulty and content based on the user's performance. After each mission, the system provides detailed performance feedback, evaluating factors such as route efficiency, safety, delivery accuracy, and completion time, helping users identify areas for improvement and refine their skills.


\section{User Stories and Use Cases}
\label{section:user-stories-use-cases}

\begin{enumerate}
    \item As a beginner drone pilot, I want to learn how to plan a drone flight by setting waypoints in QGroundControl, so that I can understand flight planning concepts.

    \item As a beginner drone pilot, I want to visualize my drone flight in a 3D environment, so that I can better understand drone behavior.
    
    \item As a beginner drone pilot, I want to simulate my planned flight using ArduPilot SITL, so that I can test my flight plan in a safe environment.
    
    \item As a beginner drone pilot, I want the system to provide training scenarios with different difficulty levels, so that I can practice progressively based on my skill level.

    \item As a beginner drone pilot, I want to receive feedback after each mission, so that I can improve my flight planning skills.

    \item As a beginner drone pilot, I want to view my performance scores on a leaderboard, so that I can track my progress and stay motivated.

\end{enumerate}

\section{Use Case Diagram}
\label{section:use-case-diagram}

\begin{figure}[H]
    \centering
    \includegraphics[width=0.6\textwidth]{examples/uml.png}
    \caption{Use Case Diagram for the Drone Training System}
    \label{fig:use-case-diagram}
\end{figure}

Use case diagram (Figure~\ref{fig:use-case-diagram}) illustrates the interactions between users and the proposed AI-driven UAV training system. The breakdown of the diagram is as follows:

\noindent \textbf{Actors:}
\begin{itemize}
    \item \textbf{Beginner Drone Pilot:} The primary user of the system.
\end{itemize}

\noindent \textbf{Use Cases:}
\begin{itemize}
    \item \textbf{Receive Mission Letter:} The user receives an AI-generated letter describing the delivery task, location, and environmental context.
    \item \textbf{Select and Pack Presents:} The user chooses gifts with weight/size properties that affect drone performance.
    \item \textbf{Plan Drone Delivery Route:} The user places waypoints on real-world terrain, sets altitude, and designs a flight path considering obstacles and safety.
    \item \textbf{Execute Drone Mission (Simulation):} The system simulates the delivery using ArduPilot SITL and MAVLink; the user can observe the flight in Unity 3D.
    \item \textbf{Receive Performance Feedback:} The user views a post-mission evaluation on route efficiency, delivery accuracy, safety, altitude control, and completion time.
    \item \textbf{Select Training Scenario / Adaptive Difficulty:} The user selects a mission or receives an AI-recommended difficulty based on their performance history.
    \item \textbf{Earn Rewards and Progress:} The user earns coins, Christmas spirit points, and unlocks harder missions.
    \item \textbf{Receive Certification:} The user receives a certificate after completing all training missions.
    \item \textbf{Enter Free Play Mode:} The user accesses continuous advanced missions after certification.
\end{itemize}

\section{User Interface Design}
\label{section:user-interface-design}
The interface supports the Christmas-themed gameplay while guiding users through the training workflow.

\begin{figure}[H]
    \centering
    \includegraphics[width=0.8\textwidth]{examples/start_page.png}
    \caption{Mockup of Start Page}
\end{figure}

\begin{figure}[H]
    \centering
    \includegraphics[width=0.8\textwidth]{examples/home_page_new.png}
    \caption{Mockup of the homepage of the application}
\end{figure}

\begin{figure}[H]
    \centering
    \includegraphics[width=0.8\textwidth]{examples/letter.png}
    \caption{Mockup of Letter Page}
\end{figure}

\begin{figure}[H]
    \centering
    \includegraphics[width=0.8\textwidth]{examples/request.png}
    \caption{Mockup of Mission Briefing}
\end{figure}

\begin{figure}[H]
    \centering
    \includegraphics[width=0.8\textwidth]{examples/playing_game1.png}
    \caption{Mockup of Simulation Interface on Unity 3D}
\end{figure}

\begin{figure}[H]
    \centering
    \includegraphics[width=0.8\textwidth]{examples/playing_game2.png}
    \caption{Mockup of Simulation Interface on Unity 3D}
\end{figure}

\begin{figure}[H]
    \centering
    \includegraphics[width=0.8\textwidth]{examples/report.png}
    \caption{Mockup of Feedback and Results Page}
\end{figure}

\begin{figure}[H]
    \centering
    \includegraphics[width=0.8\textwidth]{examples/report_2.png}
    \caption{Mockup of Feedback and Results Page}
\end{figure}

\begin{figure}[H]
    \centering
    \includegraphics[width=0.8\textwidth]{examples/report_3.png}
    \caption{Mockup of Feedback and Results Page}
\end{figure}

\begin{figure}[H]
    \centering
    \includegraphics[width=0.8\textwidth]{examples/report_4.png}
    \caption{Mockup of Feedback and Results Page}
\end{figure}

\noindent 1. Homepage:
\begin{itemize}
    \item Main entry point with game title and Christmas visual design.
    \item Options: Play, How to Play, Settings, About.
    \end{itemize}

\noindent 2. Mission Selection Page:
\begin{itemize}
    \item Displays available delivery missions based on user progress.
    \item Each mission card shows title, description, difficulty, and terrain location.
    \item Locked missions require completing previous tasks.
    \item Shows performance summary: missions completed, coins, Christmas spirit progress, average score.
\end{itemize}

\noindent 3. Mission Briefing and Letter Page:
\begin{itemize}
    \item Displays the AI-generated mission letter with the child's request, location, and environmental hints.
\end{itemize}

\noindent 4. Present Packing Page:
\begin{itemize}
    \item Select and pack presents with weight/size properties.
    \item Uses in-game coins to select gifts.
    \item Gift selection affects drone performance during the mission.
\end{itemize}

\noindent 5. Mission Planning Interface:
\begin{itemize}
    \item Real-world GIS terrain map for route planning.
    \item Place waypoints, set altitude, design flight path.
    \item Shows terrain features, obstacles, restricted areas, and delivery targets.
    \item Connected to QGroundControl workflow.
\end{itemize}

\noindent 6. Simulation Interface (Unity 3D):
\begin{itemize}
    \item 3D environment visualizing the drone delivery.
    \item Real-time drone movement, terrain, and delivery targets.
    \item Camera-follow functionality.
    \item ArduPilot SITL with MAVLink integration.
\end{itemize}

\noindent 7. Feedback and Results Page:
\begin{itemize}
    \item Post-mission performance metrics: success, route efficiency, delivery accuracy, safety violations, altitude control, completion time.
    \item Personalized feedback and improvement suggestions.
    \item Coins earned, Christmas spirit gained, rewards/unlocks.
\end{itemize}
